% Esqueleto IEEE (conferência) para o artigo da disciplina PGENE601.
% Só a estrutura: o texto, os argumentos e as conclusões são da autora.
% Compilar no Overleaf (classe IEEEtran já vem instalada) ou localmente com pdflatex.
\documentclass[conference]{IEEEtran}
\usepackage[utf8]{inputenc}
\usepackage[T1]{fontenc}
\usepackage[brazilian]{babel}
\usepackage{booktabs}
\usepackage{graphicx}
\usepackage{cite}

\begin{document}

\title{TÍTULO A DEFINIR}

\author{\IEEEauthorblockN{Fernanda Alice}
\IEEEauthorblockA{PPGINF, Universidade Federal do Amazonas\\
fernandaalice2003@gmail.com}}

\maketitle

\begin{abstract}
% Quatro a seis frases: problema, abordagem, dados, resultado principal, limite.
\end{abstract}

\begin{IEEEkeywords}
% ESBMC, Python, verificação limitada de modelos, sistemas ciberfísicos, modelos de linguagem
\end{IEEEkeywords}

\section{Introdução}
% Problema; lacuna; o que este trabalho faz; contribuições em lista curta.

\section{Fundamentação e Trabalhos Relacionados}
% ESBMC-Python; mypy; pytest; LLM para detectar bugs. Só citar o que foi conferido.

\section{Método}
\subsection{Conjunto de dados}
% 47 bugs reais de 20 projetos ciberfísicos; como foram escolhidos; versão com bug e corrigida.
\subsection{Pipeline}
% Detector (LLM), grounding, harness, especificação de tipos, ESBMC, replay em CPython.
\subsection{Linhas de comparação}
% mypy (padrão e --check-untyped-defs); testes do PR; como cada um foi medido.
\subsection{Critério de acerto}
% Gabarito pela linha do patch; níveis expressão, função, arquivo.

\section{Resultados}
% Tabela comparativa; funil 47 -> hipótese certa -> chega ao ESBMC -> confirmado.

\section{Discussão}
% Por que o funil afunila; limites do ESBMC-Python; o que cada ferramenta pega que as outras não.

\section{Ameaças à Validade}
% Mesma amostra usada para corrigir erros; gabarito pela linha; um modelo; ESBMC mutável.

\section{Conclusão}

\section*{Declaração de uso de IA}
% Declarar o uso real (código e análise de dados; revisão de texto), conforme as diretrizes.

\bibliographystyle{IEEEtran}
% \bibliography{referencias}

\end{document}
